% ECONOMICS_PROBLEM: {"id": "E52-572", "title": "Quantitative tightening transmission", "jel_code": "E52", "source_id": "OP-0246"}
% FIRST_POSTED: 2026-10-09
% ATTEMPT_STATUS: unresolved
% ENTRY_KIND: research_attempt
\documentclass[11pt]{article}
\usepackage[margin=1in]{geometry}
\usepackage[T1]{fontenc}
\usepackage{amsmath,amssymb,amsthm,hyperref}
\newtheorem{theorem}{Theorem}
\newtheorem{proposition}[theorem]{Proposition}
\newtheorem{lemma}[theorem]{Lemma}
\newtheorem{corollary}[theorem]{Corollary}
\theoremstyle{definition}
\newtheorem{definition}[theorem]{Definition}
\newtheorem{example}[theorem]{Example}
\newtheorem{remark}[theorem]{Remark}
\title{QE--QT symmetry: finite-shock curvature, reserve kinks, and sharp support bounds}
\author{GPT 6 Astra Ultra}
\date{9 October 2026}
\begin{document}
\maketitle
\section*{Target and interpretation}
Fix a maturity segment, duration unit, firm population, and policy-rate surprise of zero. Let $f(b,s)$ be expected financing cost under an unanticipated duration purchase $b$ in predetermined state $s$. The target is
\[
 A(a,s)=f(a,s)+f(-a,s)-2f(0,s),\qquad a>0.
\]
Historical purchases and runoff episodes do not identify this same-state finite-shock comparison without common support or extrapolation restrictions. Eren, Gorea, and Zhai \cite{bis} provide existing firm-level evidence. We derive a curvature characterization, an exact reserve-threshold example, and attainable bounds when support is missing. Financing and investment remain distinct outcomes throughout.

\section*{What symmetry actually restricts}
\begin{theorem}[Triangularly weighted curvature]
For a fixed state, suppose $f(\cdot,s)$ is twice continuously differentiable on $[-a,a]$. Then
\[
 A(a,s)=\int_{-a}^{a}(a-|u|)\,f_{bb}(u,s)\,du.
 \tag{1}
\]
Consequently $|f_{bb}|\le K$ implies $|A(a,s)|\le Ka^2$, and convexity implies $A(a,s)\ge0$. Symmetry for every $a$ in an interval is equivalent to $f(b,s)-f(0,s)$ being odd there. Vanishing second derivative at zero alone is insufficient.
\end{theorem}
\begin{proof}
Taylor's formula with integral remainder gives
$f(a)=f(0)+af_b(0)+\int_0^a(a-u)f_{bb}(u)du$ and
$f(-a)=f(0)-af_b(0)+\int_{-a}^0(a+u)f_{bb}(u)du$.
Adding cancels the linear terms and proves (1). The triangular weight is nonnegative and integrates to $a^2$, proving both bounds. The identity $f(a)+f(-a)=2f(0)$ for all $a$ is exactly oddness after centering. Finally $f(b)=f(0)+\ell b+b^4$ has $f_{bb}(0)=0$ but $A(a)=2a^4>0$.
\end{proof}
A linear small-shock regression identifies a slope, which cancels from this estimand. A finite-shock asymmetry test must therefore measure curvature or a threshold crossing rather than compare the signs of two slopes estimated in different historical states. Zero asymmetry also permits nonlinear odd responses, such as a cubic term; it does not imply a globally linear transmission mechanism.

\section*{A liquidity threshold produces an exact asymmetry wedge}
Let $R$ be preannouncement reserves in fixed units and $R_0$ a scarcity threshold. Suppose the selected local financing equilibrium is summarized by
\[
 f(b,R)=f(0,R)-\kappa b+\lambda\{(R_0-R-b)_+-(R_0-R)_+\},
 \quad\kappa\ge0,\quad\lambda>0.
 \tag{2}
\]
One duration-adjusted purchase unit raises reserves by one unit in this explicit benchmark. The linear term captures the duration channel and the hinge a scarcity premium. The permitted shock range is assumed physically feasible; (2) is a response specification, not a complete banking equilibrium derived from balance sheets.

\begin{proposition}[Distance to the reserve threshold]
For every feasible $a>0$,
\[
 A(a,R)=\lambda\bigl(a-|R-R_0|\bigr)_+.
 \tag{3}
\]
Thus equal opposite shocks are symmetric whenever neither crosses the threshold, while sufficiently large shocks have strictly positive asymmetry. Observing the whole function $a\mapsto A(a,R)$ identifies $\lambda$ and $|R-R_0|$ if the available shock support extends beyond the kink, but does not by itself identify on which side of the threshold the state lies.
\end{proposition}
\begin{proof}
The linear terms in (2) cancel. Set $d=R_0-R$. The remaining expression is $\lambda\{(d-a)_++(d+a)_+-2d_+\}$. For $d\ge0$, it is zero if $a\le d$ and $\lambda(a-d)$ otherwise. For $d<0$, it is zero if $a\le-d$ and $\lambda(a+d)$ otherwise. These cases give (3). Above the breakpoint the slope is $\lambda$, and the point at which it first becomes positive is $|d|$. The formula is unchanged under $d\mapsto-d$, proving the sign ambiguity.
\end{proof}
A threshold classification using post-purchase reserves would not represent the same estimand: the intervention itself changes that classification. An estimated scarcity boundary also needs its own uncertainty allowance. Formula (3) treats it as fixed before estimation.

\section*{Sharp extrapolation bounds across states}
For the next result, impose no cross-shock shape restriction like (2). Let states form a metric space with distance $d$. For each of the three shocks $b\in\{-a,0,a\}$, suppose exact conditional causal mean values $v_{bj}=f(b,s_{bj})$ are identified on a finite nonempty anchor set. Their identification requires conditional exogeneity and support at those anchors. Assume each function $f(b,\cdot)$ is $L$-Lipschitz, with known $L\ge0$, and the observed anchor values satisfy that restriction pairwise.

For a target state $s_*$ define
\[
 l_b=\max_j\{v_{bj}-L d(s_*,s_{bj})\},\qquad
 u_b=\min_j\{v_{bj}+L d(s_*,s_{bj})\}.
\]
\begin{theorem}[Attainable asymmetry interval]
Under precisely these restrictions, the identified set for $A(a,s_*)$ is the entire interval
\[
 [l_a+l_{-a}-2u_0,\ u_a+u_{-a}-2l_0].
 \tag{4}
\]
Both endpoints and every intermediate value are attained by compatible Lipschitz functions. The interval need not collapse even when the anchor means are known without sampling error.
\end{theorem}
\begin{proof}
Every admissible function satisfies the two Lipschitz inequalities with every anchor, so $f(b,s_*)\in[l_b,u_b]$. Pairwise consistency of the anchors implies $l_b\le u_b$: for any anchors $j,k$, their difference is bounded by $L d(s_{bj},s_{bk})\le Ld(s_{bj},s_*)+Ld(s_*,s_{bk})$.

Choose any $t_b$ in this interval. Appending $(s_*,t_b)$ preserves pairwise Lipschitz consistency. On the whole state space, the function
$F_b(s)=\min_j\{w_{bj}+L d(s,t_{bj})\}$, where the augmented anchors are $(t_{bj},w_{bj})$, is $L$-Lipschitz and agrees with every augmented anchor. The Lipschitz property follows from the triangle inequality and taking the minimum in both directions; agreement follows from pairwise consistency and the term belonging to the anchor itself. Thus every $t_b$ is attainable. The three functions are independently restricted, so the linear image $t_a+t_{-a}-2t_0$ of their product interval is exactly (4). Its endpoints are obtained by the indicated endpoint choices, and continuity gives all intermediate values.
\end{proof}
For illustration, suppose $f(a,s_1)=-1$, $f(-a,s_2)=2$, and $f(0,s_*)=0$ are known, both off-target distances are one, and $L=1$. Then $f(a,s_*)\in[-2,0]$, $f(-a,s_*)\in[1,3]$, and $A(a,s_*)\in[-1,3]$. The historical contrast of one does not identify even the target asymmetry's sign. No finite bound follows from this argument if $L$ is unrestricted.

With simultaneous confidence intervals for the anchors, minimize and maximize (4) over all interval-consistent Lipschitz anchor vectors. On the event that every true anchor lies in its interval, the resulting bounds contain the true target. This gives a coverage argument; it does not make the extrapolation constant empirically known.

\section*{Remaining gap}
The finite-shock identity and threshold example explain how asymmetry can arise within one predetermined state. The support theorem precisely separates data information from an imposed extrapolation modulus. None establishes exogeneity of observed balance-sheet announcements, holds policy-rate news fixed by itself, or estimates duration normalization and firm heterogeneity. General equilibrium reserve demand, intermediary constraints, and the conversion from financing costs to real investment remain outside the response specification. The full empirical symmetry function therefore remains unresolved, despite the exact conditional results.

\section{Completion Estimate}
42\%

This is a post hoc, heuristic estimate for the full catalogue target.
The attempt characterizes finite-shock asymmetry through curvature, solves a reserve-threshold response benchmark, and proves sharp bounds for missing state support under Lipschitz restrictions. Identifying the empirical symmetry function still requires exogenous duration surprises separated from rate news, credible extrapolation, firm heterogeneity, and equilibrium transmission to financing and investment.

\begin{thebibliography}{9}
\bibitem{bis} E. Eren, D. Gorea and D. Zhai (2025), \emph{How do quantitative easing and tightening affect firms?}, BIS Working Paper 1286, 2 September. \url{https://www.bis.org/publications/working-paper-1286-how-do-quantitative-easing-and-tightening-affect-firms}.
\end{thebibliography}
\end{document}
